\documentclass[11pt,a4paper]{article}
\usepackage[utf8]{inputenc}
\usepackage[T1]{fontenc}
\usepackage[margin=0.75in]{geometry}
\usepackage{enumitem}
\usepackage{titlesec}
\usepackage{xcolor}
\usepackage{hyperref}
\usepackage{lmodern}

\definecolor{primarycolor}{RGB}{33,37,41}
\definecolor{secondarycolor}{RGB}{0,92,140}
\definecolor{accentcolor}{RGB}{90,100,110}
\hypersetup{colorlinks=true,linkcolor=secondarycolor,urlcolor=secondarycolor}
\titleformat{\section}{\large\bfseries\color{primarycolor}}{}{0em}{}
\titlespacing*{\section}{0pt}{1.1em}{0.45em}
\setlist[itemize]{leftmargin=*,itemsep=0.15em,topsep=0.25em}
\pagestyle{empty}
\setlength{\parindent}{0pt}
\setlength{\parskip}{0.25em}

\begin{document}
\begin{center}
{\Huge\bfseries\color{primarycolor} KATARZYNA WIŚNIEWSKA}\\[0.3em]
{\large Pflegehelferin}\\[0.3em]
Habsburgerstraße 41, 79104 Freiburg im Breisgau | k.wisniewska@example.com | +49 761 555 0412
\end{center}

\section*{PROFIL}
Pflegehelferin mit mehrjähriger Erfahrung in der stationären Altenpflege und fünf Jahren Berufserfahrung als Krankenschwester in Polen. Sicher in der Grundpflege und Behandlungspflege nach LG1 sowie im Umgang mit Menschen mit Demenz. Ruhig, zuverlässig und geduldig; vertraut mit der Pflegedokumentation in Vivendi PD.

\section*{BERUFSERFAHRUNG}
\textbf{Pflegehelferin, Seniorenzentrum Am Lindenhof, Freiburg} \hfill \textit{04/2020–heute}
\begin{itemize}
\item Grundpflege für 10 bis 12 Bewohnerinnen und Bewohner pro Schicht: Körperpflege, Mobilisation, Lagerung und Hilfe beim Essen.
\item Behandlungspflege nach LG1: Blutdruck- und Blutzuckermessung, Kompressionsstrümpfe sowie Medikamente nach Plan reichen.
\item Pflegedokumentation in Vivendi PD.
\item Mitarbeit im Früh- und Spätdienst sowie an Wochenenden.
\item Einarbeitung von 7 neuen Pflegehelferinnen und Praktikanten in den ersten zwei Wochen nach dem Einarbeitungsplan der Einrichtung (seit 2021).
\item Leitung einer wöchentlichen Singgruppe für 6 bis 8 Bewohnerinnen und Bewohner mit Demenz (seit 2022).
\end{itemize}

\textbf{Aushilfe Verkauf, Bäckerei Schwarzwaldkorn, Freiburg} \hfill \textit{09/2019–03/2020}
\begin{itemize}
\item Verkauf, Kasse und Reinigung.
\end{itemize}

\textbf{Krankenschwester (Pielęgniarka), Wojewódzki Szpital Specjalistyczny, Opole, Polen} \hfill \textit{10/2014–06/2019}
\begin{itemize}
\item Pflege auf einer 32-Betten-Station für Innere Medizin.
\item Grund- und Behandlungspflege, Infusionen, Wundversorgung unter Aufsicht der Stationsärztin sowie Vorbereitung von Visiten.
\end{itemize}

\section*{AUSGEWÄHLTES PROJEKT}
\textbf{Sturzprophylaxe im Wohnbereich 2, Seniorenzentrum Am Lindenhof} \hfill \textit{2023}
\begin{itemize}
\item Gemeinsam mit der Wohnbereichsleitung Sturzprotokolle aus sechs Monaten ausgewertet und für 9 Bewohnerinnen und Bewohner Nachtlichter und feste Toilettengänge eingeführt.
\item Die Zahl der Stürze im Wohnbereich sank im folgenden Halbjahr von 14 auf 8.
\end{itemize}

\section*{AUSBILDUNG UND ANERKENNUNG}
Bachelorabschluss in Krankenpflege (Pielęgniarstwo), Medizinische Universität Opole | 2014\\
Anerkennungsverfahren als Pflegefachfrau beim Regierungspräsidium Stuttgart: Defizitbescheid erhalten; Kenntnisprüfung für März 2027 vorgesehen.

\section*{WEITERBILDUNG}
Behandlungspflege LG1 (Basiskurs) | 2021\\
Expertenstandard Sturzprophylaxe (interne Schulung) | 2023\\
Erste-Hilfe-Kurs | 2024

\section*{KENNTNISSE}
Grundpflege und Mobilisation | Behandlungspflege nach LG1 | Umgang mit Menschen mit Demenz | Pflegedokumentation: Vivendi PD | Einarbeitung von Pflegehelferinnen und Praktikanten

\section*{SPRACHEN}
Polnisch (Muttersprache) | Deutsch (B2, telc Deutsch B1-B2 Pflege; Zertifikat 2021) | Englisch (Grundkenntnisse)
\end{document}